% !TeX root = ../../tesis.tex

\subsection{Transport-gis: Su arquitectura y funcionamiento}
\label{subsec:transport-gis-arquitectura-funcionamiento}

\textbf{Transport-gis} se organiza en cinco módulos funcionales
interdependientes: un módulo de \textbf{análisis} (\texttt{analysis/}), que
gestiona la consulta a los servicios externos y el procesamiento de indicadores;
un módulo de \textbf{datos} (\texttt{data/processed/}), que almacena las capas
geoespaciales en formato \termino{GeoPackage} (\texttt{.gpkg}) bajo control de
Git~\textsc{lfs}; un módulo \textbf{cartográfico} (\texttt{maps/}), que agrupa
los proyectos \textsc{qgis}, las plantillas tipográficas y los archivos de
exportación; un módulo de \textbf{conectividad analítica} (\texttt{tableau/}),
que contiene un \termino{Web Data Connector} para \termino{Tableau Desktop}, un
servidor proxy para geometrías complejas y utilidades Python de exportación; y
un módulo de \textbf{orquestación} (\texttt{pipeline/}), que coordina la
ejecución secuencial del proceso completo. La Figura~\ref{fig:tg-pipeline}
ilustra el flujo de datos entre estos módulos y los componentes externos de la
arquitectura de la investigación.

% =============================================================================
\subsubsection{Flujo de producción de datos}
\label{subsubsec:tg-flujo-datos}
% =============================================================================

El proceso de producción cartográfica parte de dos fuentes primarias externas:
los \textit{feeds} \gtfs{} de la \textsc{semovi} (2024), que contienen los
horarios, rutas y paradas de los sistemas de transporte de la \cdmx{}; y la
cartografía base del \textsc{inegi} (2023), que provee los límites
político-administrativos y la infraestructura vial de la \zmvm{}. Estas fuentes
alimentan a \textbf{Apimetro} —que las transforma en una \api{} \textsc{rest}
consumible con la topología del grafo de la red— y a \textbf{VFTModel}, que las
procesa para calcular los indicadores de eficiencia nodal.

Dentro del repositorio de \textbf{Transport-gis}, el script
\texttt{apimetro\_client.py} consulta la \api{} de \textbf{Apimetro} para
recuperar la topología del grafo y los atributos de cada modo de transporte,
mientras que \texttt{vft\_runner.py} ejecuta las rutinas del \textbf{VFTModel}
para producir los indicadores por nodo en formato \texttt{.csv}. El orquestador
\texttt{pipeline/run\_pipeline.py} gestiona la secuencia completa de
transformaciones hasta exportar las capas procesadas al directorio
\texttt{data/processed/}. Los fragmentos de código de estos componentes se
documentan en el Anexo~\ref{anx:transportgis-tecnico}.

\begin{figure}[htbp]
\centering
% Diagrama: Pipeline de producción cartográfica de Transport-gis
% Usado en: 4-Capas-Codigo/4-3-Visualizacion.tex  (fig:tg-pipeline)
\begin{tikzpicture}[
  srcbox/.style = {rectangle, rounded corners=3pt, draw=gray!55, fill=gray!8,
                   minimum width=3.0cm, minimum height=0.90cm,
                   align=center, font=\small},
  extbox/.style = {rectangle, rounded corners=3pt, draw=unamAzul, fill=unamAzul!10,
                   minimum width=3.2cm, minimum height=0.90cm,
                   align=center, font=\small},
  intbox/.style = {rectangle, rounded corners=3pt, draw=unamAzul!60, fill=unamGrisClaro,
                   minimum width=6.4cm, minimum height=1.05cm,
                   align=center, font=\small},
  outbox/.style = {rectangle, rounded corners=3pt, draw=unamOro, fill=unamOro!15,
                   minimum width=3.2cm, minimum height=0.90cm,
                   align=center, font=\small},
  ->, >=Stealth, thick
]
%% — Fuentes externas —
\node[srcbox] (gtfs)  at (0.0, 9.2) {\gtfs{} SEMOVI~2024};
\node[srcbox] (inegi) at (5.4, 9.2) {INEGI~2023 (cartografía)};

%% — Componentes externos —
\node[extbox] (api) at (0.0, 7.5) {\textbf{Apimetro}\\{\footnotesize API REST (Go)}};
\node[extbox] (vft) at (5.4, 7.5) {\textbf{VFTModel}\\{\footnotesize (Python / NetworkX)}};

%% — Transport-gis: análisis —
\node[intbox] (scripts) at (2.7, 5.5)
  {{\footnotesize\texttt{apimetro\_client.py} $\cdot$ \texttt{vft\_runner.py}}\\
   {\footnotesize\texttt{pipeline/run\_pipeline.py}}};

%% — Transport-gis: datos —
\node[intbox] (gpkg) at (2.7, 4.0)
  {{\small\texttt{data/processed/}}\\
   {\footnotesize 10 capas \texttt{.gpkg} (Git~LFS)}};

%% — Transport-gis: QGIS —
\node[intbox] (qgis) at (2.7, 2.5)
  {{\small QGIS~3.x — \texttt{maps/projects/*.qgz}}\\
   {\footnotesize\texttt{maps/templates/*.qpt} (paletas UNAM)}};

%% — Transport-gis: exports —
\node[intbox] (exports) at (2.7, 1.0)
  {{\small\texttt{maps/exports/}}\\
   {\footnotesize\texttt{*.pdf} $\cdot$ \texttt{*.png}}};

%% — Destino final —
\node[outbox] (tesis) at (2.7, -0.5)
  {\textbf{Tesis LaTeX}\\{\footnotesize\texttt{Figures/}}};

%% — Flechas: fuentes → externos —
\draw[color=gray!60] (gtfs)  -- (api);
\draw[color=gray!60] (inegi) -- (vft);

%% — Flechas: externos → scripts (convergentes, con codo exterior) —
\draw[color=unamAzul!70] (api.south) -- ++(0,-0.55) -| (scripts.north west);
\draw[color=unamAzul!70] (vft.south) -- ++(0,-0.55) -| (scripts.north east);

%% — Flujo interno —
\draw[color=unamAzul!70] (scripts) -- (gpkg);
\draw[color=unamAzul!70] (gpkg)    -- (qgis);
\draw[color=unamAzul!70] (qgis)    -- (exports);

%% — Exports → LaTeX —
\draw[color=unamOro] (exports) -- (tesis);

%% — Recuadro del repositorio Transport-gis (en fondo) —
\begin{scope}[on background layer]
  \node[draw=unamAzul!40, dashed, rounded corners=6pt, fill=unamAzul!3,
        fit={(scripts)(gpkg)(qgis)(exports)},
        inner sep=14pt] (tgbox) {};
\end{scope}
%% — Label DENTRO del recuadro, esquina sup-izq —
\node[font=\small\bfseries\color{unamAzul}, anchor=north west]
  at ($(tgbox.north west) + (4pt,-3pt)$) {Transport-gis-zmvm-mjg};

\end{tikzpicture}

\caption[Pipeline de producción cartográfica de Transport-gis]%
  {Flujo de datos del sistema \textbf{Transport-gis-zmvm-mjg}.
Los componentes externos \textbf{Apimetro} y \textbf{VFTModel} proveen la
información primaria;
dentro del repositorio, los \textit{scripts} de análisis consolidan los datos en
capas \texttt{.gpkg}
  que QGIS transforma en los mapas de presentación exportados a la tesis.}
\label{fig:tg-pipeline}
\end{figure}

% =============================================================================
\subsubsection{Integración con Tableau Desktop}
\label{subsubsec:tg-tableau}
% =============================================================================

Junto al flujo cartográfico de \textsc{qgis}, \textbf{Transport-gis} cuenta con
herramientas para alimentar tableros en \termino{Tableau Desktop}. La distinción
entre ambas salidas responde a necesidades distintas: \textsc{qgis} produce
representaciones estáticas de alta resolución para la tesis impresa; Tableau
habilita la exploración temporal por sistema, alcaldía y período sin requerir
reprocesar las capas ya calculadas.

El módulo de conectividad analítica articula dos mecanismos según la naturaleza
del dato. La afluencia mensual y el catálogo de estaciones son estructuras
tabulares; para ellas se emplea un \termino{Web Data Connector} (\textsc{wdc}):
una página \textsc{html} con JavaScript que invoca los endpoints analíticos de
\textbf{Apimetro} y mapea las respuestas \textsc{json} a columnas tipadas
directamente en Tableau. Los trazados y polígonos de cobertura, en cambio, son
geometrías complejas (\texttt{MultiLineString}, \texttt{MultiPolygon}) que el
\textsc{wdc} no puede transmitir; para estas capas se utiliza el conector
\termino{Spatial File} de Tableau, que acepta \textsc{url}s \textsc{http} que
sirvan \geojson{}. El servidor \textbf{GeoProxy} (FastAPI, puerto~\texttt{5050})
cubre ese rol: re-expone como rutas locales lo que \textbf{Apimetro} y
\textbf{VFTModel} ya sirven en sus propios puertos, sin transformar el
contenido.

La implementación completa de ambos mecanismos —código \textsc{wdc}, servidor
\textbf{GeoProxy} y utilidades de exportación— se documenta en el Anexo~
\ref{anx:transportgis-tecnico}.
